\documentclass[10pt,a4paper]{article}
\usepackage[utf8]{inputenc}
\usepackage[a4paper,left=3cm,right=2cm]{geometry}
\usepackage{amsmath}
\usepackage[thmmarks, amsmath, thref]{ntheorem}
\usepackage{amsfonts}
\usepackage{amssymb}
\usepackage{fancyhdr}
\pagestyle{fancy}
\usepackage{anysize}
\usepackage{graphicx}
\usepackage{float}
\usepackage{hyperref}
\usepackage{enumitem}
\usepackage{verbatim}
\newtheorem{theorem}{Theorem}
\theoremsymbol{\ensuremath{\blacksquare}}
\newtheorem*{solution}{Solución:}

\usepackage{xcolor}

% Margenes y formalidades

\textwidth 6.25in % Margen Derecho y ancho del texto
\textheight 8.5in
\oddsidemargin 0in % Margen Izquierdo
\headheight 0in % Largo del texto
\leftmargin 1in
\parindent 0em % Salto entre la indentación (Sangría)https://www.overleaf.com/4519127669sdpptkhvsbfn
\topmargin -0.5in % Margen Superior
\parskip 0.5ex % Salto entre parrafos
\baselineskip 1pt



\lhead{
\includegraphics[scale=0.08]{Figuras/logoUSM.png}} % texto izquierda de la cabecera
%\chead{}                                      % texto centro de la cabecera
\rhead{Universidad Técnica Federico Santa María\\Ingeniería Comercial } % número de página a la derecha
\renewcommand{\headrulewidth}{0.4pt} % grosor de la línea de la cabecera
\renewcommand{\footrulewidth}{0.4pt} % grosor de la línea del pie

\begin{document}
\vspace{5cm}



\begin{center}
\begin{large}
\textbf{\\CERTAMEN 2:}
\textbf{Econometría ICS-294}
\end{large}\\
\vspace{0.2cm}
I Semestre 2025\\
Profesores: Jocelyn Tapia Stefanoni, Franco Antonucci Fernández y Ricardo Pruneda Urz\'ua\\
%Ayudante: 
\end{center}

\textbf{Instrucciones:} 

\begin{itemize}

    \item Usted dispone de 90 minutos  para intentar los 100 puntos del certamen.
    \item Cuide la presentación y redacción de sus preguntas.
    \item No puede utilizar celular, computador o tablet, pero sí una calculadora TI.
     
\end{itemize}

\section{RdA3.2  Preguntas cortas (30 puntos)}
{\bf Elija solo 3 preguntas y responda. }
\begin{enumerate}
   \item   (10 puntos) 
   Considere el siguiente modelo de regresión poblacional para el logaritmo del salario:
    
    $$\ln{wage} = \beta_0 + \beta_1 educ + \beta_2 exper + \mu$$
    
    Donde $educ$: educación; $exper$: años de experiencia. Al estimar el modelo se obtienen los siguientes valores:
    \begin{table}[H] 
    \centering
    \caption{Resultados de la regresión}
    \begin{tabular}{lcc}
    \hline
    \hline
    \textbf{Variable} & \textbf{Coeficiente} & \textbf{Error Estándar} \\
    \hline
    Intercepto & 1.20$^{*}$ & 0.15 \\
    educ & 0.08$^{***}$ & 0.01 \\
    exper & 0.03$^{***}$ & 0.005 \\
    \hline
    \multicolumn{3}{l}{\textbf{n} = 1000} \\
    \multicolumn{3}{l}{$R^2 = 0.32$} \\
    \multicolumn{3}{l}{$F = 235.4$ (gl = 2,997)} \\
    \hline
    \hline
    \textit{Nota:} $^{*}p<0.1$; $^{**}p<0.05$; $^{***}p<0.01$
    \end{tabular}
    \end{table}
    
    {Con lo cual un investigador le dice a usted que no hay retornos positivos por más experiencia al 5\% de significancia. Comente.}

     \item (10 puntos) ¿Cómo podría reducir la varianza de los estimadores MCO en un estudio econométrico real? Señale al menos dos formas

      \item 
    (10 puntos) Un estimador es consistente cuando $n \longrightarrow \infty$; en este caso la distribución del estimador colapsa hasta el valor del parámetro que se está estimando. Esto significa que el parámetro es consistente e insesgado. Comente.

      \item (10 puntos)  En un modelo de regresión lineal simple estimado usando una base de datos con  150 observaciones, la suma de residuos al cuadrado es 19.000, determine  $\hat{\sigma}^2$.
    \end{enumerate}

    \section{RdA3.2  Ejercicios (60 puntos)}


{\bf A continuación elija solo 3 preguntas y responda. }

\begin{enumerate}
 \item  Una investigación se encuentra analizando los principales determinantes del peso de los recién nacidos.
Se dispone de una base de datos con las siguientes variables: 

\begin{itemize}
    \item 
bwght: peso al nacer, en libras.
\item cigs: número promedio de cigarros diarios que fumó la madre durante el embarazo.
\item parity: orden de nacimiento del niño.
\item faminc: ingreso anual de la familia.
\item mo\_educ: años de escolaridad de la madre.
\item fa\_educ: años de escolaridad del padre.
\item male: variable dicotómica que toma el valor 1 si el recién nacido es varón.
\end{itemize}

Los resultados obtenidos están representados en la siguiente ecuación, donde en paréntesis aparecen los errores estándar de los estimadores.
    
{\scriptsize
\begin{equation*}
\begin{array}{ll}
\widehat{log(bwght)}  = & 1,1 -0,0052 cigs+0,011 log( faminc) +0,017 parity+0,034 male-0,0065mo\_educ +0,005fa\_educ\\
& (0,38) \quad (0, 001) \quad \quad \quad \quad (0,004)\quad \quad \quad \quad \quad (0,006) \quad \quad \quad \quad (0,011)\quad \quad \quad \quad  (0,03)\quad \quad \quad \quad \quad (0,0046)\\
 & n =1.191 \quad R^2=  0,493
\end{array}
\end{equation*}}


\begin{enumerate}
    
    \item (5 puntos) Analice el efecto de un aumento del 10\% en el ingreso anual de la familia sobre el  peso al nacer. Interprete este efecto en términos económicos y evalúe su significancia estadística utilizando un nivel de significancia del 1\% (\( \alpha = 0{,}01 \)). ¿Se puede concluir que el ingreso familiar tiene un efecto significativo sobre el peso al nacer?

    \item (5 puntos) Evalúe la significancia estadística individual de las variables \textit{escolaridad del padre} (\( \text{fatheduc} \)) y \textit{escolaridad de la madre} (\( \text{motheduc} \)) en el modelo estimado. Para ello, realice las pruebas t correspondientes y determine si alguna de estas variables tiene un efecto significativo sobre el peso al nacer, utilizando un nivel de significancia del 5\%.

    \item \textbf{(10 puntos)} Se presenta a continuación un modelo alternativo estimado para explicar el logaritmo del peso al nacer:

    \begin{equation*}
    \begin{aligned}
    \widehat{\log(\text{bwght})} =\ & 1{,}66 - 0{,}0044\, \text{cigs} + 0{,}0093\, \log(\text{faminc}) + 0{,}016\, \text{parity} + 0{,}027\, \text{male} \\
    & (0{,}22) \quad (0{,}0009) \quad \quad (0{,}0059) \quad \quad (0{,}006) \quad \quad (0{,}010) \\
    & n = 1.191 \quad \quad R^2 = 0{,}472
    \end{aligned}
    \end{equation*}

    Utilizando información de este modelo y del modelo original que incluía las variables \textit{motheduc} y \textit{fatheduc}, realice una prueba de significancia conjunta (prueba F) para evaluar si ambas variables pueden ser excluidas del modelo sin pérdida significativa de capacidad explicativa. ¿Debe la investigación preferir el modelo completo o el modelo reducido? Justifique su elección en función de criterios estadísticos y del objetivo del análisis.

\end{enumerate}

  \item Considere la siguiente información:
 \[
n=600 \quad \hat{\beta} =  \begin{bmatrix}
1,5 \\
2 \\
3,55 
\end{bmatrix} \quad \hat{\sigma}^2 = 141925; \quad 
(X^\top X)^{-1} = 
\begin{bmatrix}
0.05987 & -0.00177 & -0.00035 \\
-0.00177 & 0.00030 & -0.00002 \\
-0.00035 & -0.00002 & 0.00001
\end{bmatrix}
\]

\begin{enumerate}
    \item (9 puntos) Determine los errores estándar de $\hat{\beta}_0$, $\hat{\beta}_1$ y $\hat{\beta}_2$.
\item (11 puntos) Plantee la hip\'otesis  que testea $\beta_2$ es el doble de  $\beta_1$. 
¿Se puede rechazar la hipótesis nula con la información que disponemos con una significancia del 1\%?
    \end{enumerate}
      \item     Suponga que, debido al tornado que hubo en Puerto Varas (Región de Los Lagos, Chile) el día Domingo 25 de Mayo de 2025, le piden a usted estimar un modelo de regresión lineal  para la fuerza de un tornado (medido en la escala Fujita). Usted toma muchos datos de tornados y estima que el modelo es:

    $$F_i = \gamma_0 + \gamma_1 \cdot V_i + \gamma_2 \cdot P_i + \gamma_3 \cdot T_i + u_i$$

    Donde $F$ es la fuerza del tornado; $V$ es la velocidad del viento ($km/h$); $P$: nivel de precipitaciones al momento de ocurrir el tornado (en mm); $T$: temperatura al momento de ocurrir el tornado (en ºC).\\

    Se tiene lo siguiente:

\begin{table}[H] \centering 
  %\caption{Regresión múltiple sobre la fuerza del tornado} 
  \label{tab:tornado_reg} 
\begin{tabular}{@{\extracolsep{5pt}}lcc} 
\\[-1.8ex]\hline 
\hline \\[-1.8ex] 
 & \multicolumn{2}{c}{\textit{Dependent variable:}} \\ 
\cline{2-3} 
\\[-1.8ex] & \multicolumn{2}{c}{Fuerza del Tornado (F)} \\ 
\\[-1.8ex] & (1) & (2) \\ 
\hline \\[-1.8ex] 
 \textit{Velocidad del Viento} & 0.025$^{***}$ & 0.020$^{***}$ \\ 
  & (0.0016) & (0.0025) \\ 
  & & \\ 
 \textit{Precipitaciones} &  & 0.030$^{***}$ \\ 
  &  & (0.0050) \\ 
  & & \\ 
 \textit{Temperatura} &  & 0.050 \\ 
  &  & (0.0300) \\ 
  & & \\ 
 Intercepto & -2.000$^{***}$ & -1.500$^{***}$ \\ 
  & (0.289) & (0.390) \\ 
  & & \\ 
\hline \\[-1.8ex] 
Observaciones & 182 & 182 \\ 
R$^{2}$ & 0.45 & 0.62 \\ 
Adjusted R$^{2}$ & 0.44 & 0.61 \\ 
Error estándar residual & 0.45 (df = 180) & 0.35 (df = 178) \\ 
F Statistic & 148.0$^{***}$ (df = 1; 180) & 96.3$^{***}$ (df = 3; 178) \\ 
\hline 
\hline \\[-1.8ex] 
\textit{Nota:}  & \multicolumn{2}{r}{$^{*}$p$<$0.1; $^{**}$p$<$0.05; $^{***}$p$<$0.01} \\
\textit{}  & \multicolumn{2}{r}{Valores en paréntesis corresponden al error estándar}
\end{tabular} 
\end{table}  

\begin{enumerate}
    \item {(5 puntos)} En términos de los parámetros del modelo, establezca la hipótesis nula de que, controlando por \textit{Precipitaciones} y \textit{Temperatura}, la \textit{Velocidad} no tiene efecto en la \textit{Fuerza del Tornado} y luego comente la aseveración de lo anterior.
     \item {(5 puntos)} Interprete los coeficientes en la regresión (2). En particular \textit{Velocidad del Viento}, ¿es mayor o menor que el coeficiente correspondiente en la regresión simple (1)? Explique.
  \item {(5 puntos)} Testee la significancia global del modelo al 1\% y explique.\footnote{Hint: considere que el valor crítico de F al 1\% con grados de libertad (3,178) es aproximadamente 3.87}
   \item {(5 puntos)} Para explicar el Fuerza del Tornado, ¿es mejor el modelo restringido o el no restringido? Argumente.
        \end{enumerate}

        \item Un grupo de analistas se encuentra realizando una investigación sobre  el mercados de propiedades de la ciudad de Boston.  
Cuentan con los siguientes datos:

\begin{itemize}
    \item Precio de las propiedades (lprice en logaritmos).
    \item La cantidad de baños de las propiedades (baths). 
    \item La cantidad de habitaciones (rooms).
        \item Antiguedad de la propiedad (age en cantidad de años y agesq los años al cuadrado). 
\end{itemize}

Sus resultados se resumen en la siguiente  tabla:

% Table created by stargazer v.5.2.3 by Marek Hlavac, Social Policy Institute. E-mail: marek.hlavac at gmail.com
% Date and time: Thu, May 29, 2025 - 09:54:11 PM
\begin{table}[!htbp] \centering 
  \caption{} 
  \label{} 
\begin{tabular}{@{\extracolsep{5pt}}lcc} 
\\[-1.8ex]\hline 
\hline \\[-1.8ex] 
 & \multicolumn{2}{c}{\textit{Dependent variable:}} \\ 
\cline{2-3} 
\\[-1.8ex] & \multicolumn{2}{c}{lprice} \\ 
\\[-1.8ex] & (1) & (2)\\ 
\hline \\[-1.8ex] 
 larea & 0.437$^{***}$ & 0.545$^{***}$ \\ 
  & (0.073) & (0.069) \\ 
  & & \\ 
 baths & 0.224$^{***}$ & 0.070$^{*}$ \\ 
  & (0.034) & (0.038) \\ 
  & & \\ 
 rooms & 0.031 & 0.041$^{*}$ \\ 
  & (0.024) & (0.023) \\ 
  & & \\ 
 age &  & $-$0.009$^{***}$ \\ 
  &  & (0.002) \\ 
  & & \\ 
 agesq &  & 0.00003$^{***}$ \\ 
  &  & (0.00001) \\ 
  & & \\ 
 Constant & 7.328$^{***}$ & 6.914$^{***}$ \\ 
  & (0.490) & (0.461) \\ 
  & & \\ 
\hline \\[-1.8ex] 
Observations & 321 & 321 \\ 
R$^{2}$ &  & 0.597 \\ 
Adjusted R$^{2}$ &  & 0.590 \\ 
Residual Std. Error & 0.305 (df = 317) & 0.280 (df = 315) \\ 
F Statistic & 114.881$^{***}$ (df = 3; 317) & 93.213$^{***}$ (df = 5; 315) \\ 
\hline 
\hline \\[-1.8ex] 
\textit{Note:}  & \multicolumn{2}{r}{$^{*}$p$<$0.1; $^{**}$p$<$0.05; $^{***}$p$<$0.01} \\ 
\end{tabular} 
\end{table} 

\begin{enumerate}
    \item (6 puntos) En base a la información de la tabla determine el $R^2$ de la primera regresión.
    \item (6 puntos) Según la segunda regresión, si el área aumenta en 1\%, cuál es el cambio esperado en el precio de la propiedad? ¿Es esta variable explicativa estadísticamente significativa? Explique.  
    \item (8 puntos) Describa la relación estimada entre el precio de las propiedades y la antig\"uedad de ésta. Cree usted que es apropiada la inclusión de la antig\"uedad en el modelo. Justifique su respuesta.
\end{enumerate}
    \end{enumerate}

    \newpage


\begin{center}
    {\Large Formulario  }

 
\end{center}

\begin{itemize}

% \item  
% {\color{red}
% $$\hat{\beta}_0 = \bar{y} - \hat{\beta}_1 \bar{x} $$
% $$\hat{\beta}_1 = \frac{\sum_{i=1}^{n}(x_i - \bar{x})(y_i-\bar{y})}{\sum_{i=1}^{n}(x_i - \bar{x})^2}$$}

\item  {\color{red} \[ \hat{\beta} = (X^{\top}X)^{-1}X^{\top}y \]}

\item 
	{\color{red}	\[ SST = \sum_{i=1}^{n} (y_i - \bar{y})^2 \quad SSE = \sum_{i=1}^{n} (\hat{y}_i - \bar{y})^2  \quad SSR = \sum_{i=1}^{n} \hat{u}_i^2\]}
		

    \item {\color{red} \[ R^2  = 1 - \frac{SSR}{SST} \quad \quad \bar{R}^2 =  1 - \frac{(n-1)SSR}{(n-k-1)SST} \]}

 \item {\color{red} \[ t_{\hat{\beta_j}}=\frac{\hat{\beta}_j - \beta_j}{\text{s.e.}(\hat{\beta}_j)} \sim t_{n-k-1} \]}

 \item {\color{red}\[
    F = \frac{SSR_{r} - SSR_{ur}}{SSR_{ur}} \cdot \frac{n - k - 1}{j} \sim F(j, n - k - 1)
    \]}

% \item {\color{red}Considere que el valor crítico a un nivel de significancia del 5\% es 1.645}
% \item {\color{red}Considere los siguientes valores críticos para t: $\alpha = 0.01 \implies t \approx 2.62$; $\alpha = 0.05 \implies t \approx 1,97$;$\alpha = 0.01 \implies t \approx 2.62$; $\alpha = 0.1 \implies t \approx 1,645$}
\item {\color{red}Considere que el valor crítico de F al 1\% con grados de libertad (2,178) es aproximadamente 4.72}

\end{itemize}

\begin{figure}[H]
    \centering
    \includegraphics[width=0.95\linewidth]{Ejercicios//ejercicios_inferencia/cert2_F2_1184.png}
\end{figure}

\begin{figure}[H]
    \centering
    \includegraphics[width=0.95\linewidth]{certamenes/normal.png}
\end{figure}


\begin{figure}
    \centering
    \includegraphics[width=0.95\linewidth]{certamenes/normal_1cola.png}
\end{figure}
    \end{document}